\section{Far-annulus crowding below the quadratic norm}
\label{crowding:section}

The far-collision hypothesis in Proposition~\ref{angselect:farcriterion}
is stronger than needed. A fixed positive fraction of pairs with bounded
far-annulus crowding suffices. We prove an exact characterization of the
absolutely continuous mass of the selected law, and a logarithmic criterion
for its full absolute continuity. These are measure-theoretic reductions,
not new dimension estimates. In particular, the characterization concerns
the selected law; it is not a necessary condition for its distance support
to have positive length.

Retain the separated probabilities, fixed angular threshold, and exponent
\(1/(s+1)<\alpha<1\) of Section~\ref{angselect:section}. Write
\[
 d\Pi_A(x,y)=a_A(x,y)\,d\mu(x)d\nu(y),\qquad
 \sigma_y=(d_y)_*(a_A(\cdot,y)\mu),\qquad
 \Lambda_A(dy,dt)=d\nu(y)\,\sigma_y(dt).
\]
Its mass \(m_A\) is positive and at most one. For \(0<\epsilon<1\), set
\begin{align*}
 H_\epsilon(x,y)
 &=\epsilon^{-1}\sigma_y
   ([|x-y|-\epsilon,|x-y|+\epsilon]),\\
 N_\epsilon(x,y)
 &=\epsilon^{-1}\int a_A(x',y)
   \mathbf1_{\{||x-y|-|x'-y||\le\epsilon\}}
   \mathbf1_{\{|x-x'|\le\epsilon^\alpha\}}\,d\mu(x'),\\
 F_\epsilon(x,y)&=H_\epsilon(x,y)-N_\epsilon(x,y).
\end{align*}
Thus \(F_\epsilon\) counts only the far source points. All these
functions are nonnegative and measurable. Proposition~\ref{angselect:nearprop}
gives
\begin{equation}\label{crowding:near}
 \int N_\epsilon\,d\Pi_A\le C\epsilon^\kappa\longrightarrow0,
 \qquad
 \kappa=\min\{\alpha(s+1)-1,\alpha(s-1)\}>0.
\end{equation}

\subsection{Retain a positive mass with bounded far crowding}

\begin{proposition}\label{crowding:positive}
Suppose \(\epsilon_j\to0\) and, for some finite \(T>0\),
\[
 m:=\limsup_{j\to\infty}
       \Pi_A\{F_{\epsilon_j}\le T\}>0.
\]
Then \(\Lambda_A\) dominates a measure \(g(y,t)\,d\nu(y)dt\) with
\[
 g\ge0,\qquad \int g\,d\nu\,dt=m,\qquad
 \int g^2\,d\nu\,dt\le Tm.
\]
In particular there is a set of pins of \(\nu\)-measure at least
\(m/4\) such that
\[
 |\Delta_y(\supp\mu)|\ge \frac{m^2}{16T}.
\]
No nesting or convergence of the pair selections is required.
\end{proposition}

\begin{proof}
Pass to a subsequence on which the displayed pair masses converge to
\(m\). Put \(G_j=\{F_{\epsilon_j}\le T\}\), and retain the positive
source kernel
\[
 \eta_{j,y}(dx)=a_A(x,y)\mathbf1_{G_j}(x,y)\,d\mu(x).
\]
Use the box mollifier
\(b_\epsilon=\epsilon^{-1}\mathbf1_{[-\epsilon/2,\epsilon/2]}\)
and set \(f_j(y,t)=b_{\epsilon_j}*(d_y)_*\eta_{j,y}(t)\).
Its correlation kernel satisfies the exact inequality
\[
 (b_\epsilon*\widetilde b_\epsilon)(u)
 =\frac{(\epsilon-|u|)_+}{\epsilon^2}
 \le\epsilon^{-1}\mathbf1_{\{|u|\le\epsilon\}}.
\]
Tonelli and positivity, dropping only the second factor
\(\mathbf1_{G_j}(x',y)\), therefore give
\begin{align*}
 \|f_j\|_{L^2(d\nu\,dt)}^2
 &\le\int_{G_j}H_{\epsilon_j}(x,y)\,d\Pi_A(x,y)\\
 &\le T\Pi_A(G_j)+\int N_{\epsilon_j}\,d\Pi_A
 \le T\Pi_A(G_j)+C\epsilon_j^\kappa.
\end{align*}
The near term is integrated over all pairs, so no additional near-crowding
deletion is needed.

All radial supports lie in one fixed bounded interval. Weak compactness in
the joint \(L^2\) space gives a limit \(g\ge0\) of mass \(m\), with
\(\|g\|_2^2\le Tm\). The selected mollifications are dominated by
the full mollifications of \(\Lambda_A\). Testing against nonnegative
continuous functions and taking limits gives
\(g\,d\nu\,dt\le\Lambda_A\), exactly as in
Proposition~\ref{selection:compactness}.

Disintegration gives \(g(y,t)dt\le\sigma_y\) for almost every pin.
Let \(h(y)=\int g(y,t)dt\), so \(0\le h\le1\) and \(\int h\,d\nu=m\).
The set \(\{h\ge m/2\}\) has measure at least \(m/2\).
Markov's inequality bounds the measure of
\(\{\int g(y,t)^2dt>4T\}\) by \(m/4\).
Their difference therefore has measure at least \(m/4\), and on it
Cauchy--Schwarz gives
\[
 |\Delta_y(\supp\mu)|
 \ge\frac{h(y)^2}{\int g(y,t)^2dt}
 \ge\frac{m^2}{16T}.
\]
The domination by \(\sigma_y\) justifies the distance support in this
inequality.
\end{proof}

\subsection{The exact amount of absolutely continuous mass}

Let \(m_{\rm ac}\) be the mass of the absolutely continuous part of
\(\Lambda_A\) relative to \(d\nu(y)dt\).

\begin{proposition}\label{crowding:mass}
For every sequence \(\epsilon_j\to0\),
\begin{equation}\label{crowding:massformula}
 m_{\rm ac}
 =\sup_{T>0}\liminf_{j\to\infty}\Pi_A\{F_{\epsilon_j}\le T\}
 =\sup_{T>0}\limsup_{j\to\infty}\Pi_A\{F_{\epsilon_j}\le T\}.
\end{equation}
Consequently a fixed positive mass at a fixed finite crowding threshold is
both sufficient and necessary for \(\Lambda_A\) to have a nonzero
absolutely continuous part. Full absolute continuity follows if
\[
 \lim_{T\to\infty}\limsup_{j\to\infty}
       \Pi_A\{F_{\epsilon_j}>T\}=0.
\]
\end{proposition}

\begin{proof}
Proposition~\ref{crowding:positive} shows that, for each finite \(T\),
the limsup in \eqref{crowding:massformula} is at most \(m_{\rm ac}\).
This is immediate also when that limsup is zero.

For the reverse inequality write the absolutely continuous part as
\(f(y,t)\,d\nu(y)dt\). The Lebesgue decomposition disintegrates:
for \(\nu\)-almost every \(y\), the absolutely continuous part of
\(\sigma_y\) is \(f(y,t)dt\). Indeed a joint null set carrying the
singular remainder has Lebesgue-null fibres for almost every pin, by
Fubini; disintegrating the two positive parts then proves the assertion.
The differentiation theorem for a finite measure on the line gives
\[
 \frac{\sigma_y([t-\epsilon,t+\epsilon])}{\epsilon}
       \longrightarrow 2f(y,t)
\]
at Lebesgue-almost every \(t\), hence also at
\(f(y,t)\,d\nu(y)dt\)-almost every \((y,t)\).
Since \(F_\epsilon\le H_\epsilon\), pushforward and Fatou give
\[
 \liminf_j\Pi_A\{F_{\epsilon_j}\le T\}
 \ge\int_{\{2f<T\}} f(y,t)\,d\nu(y)dt.
\]
Let \(T\to\infty\). The integral tends to \(m_{\rm ac}\), proving
both equalities. The final assertion follows by taking complements in
the finite measure \(\Pi_A\), whose mass equals that of \(\Lambda_A\).
\end{proof}

Thus a bound for the whole first moment \(\int F_\epsilon\,d\Pi_A\)
is unnecessary. For positive length it is enough to bound crowding on a
fixed positive mass; to obtain full absolute continuity it is enough that
the far-crowding distributions are uniformly tight along one sequence.
Neither statement estimates these quantities using \(d,D\).

\subsection{A logarithmic sufficient estimate}

\begin{proposition}\label{crowding:entropy}
If, for some sequence \(\epsilon_j\to0\),
\begin{equation}\label{crowding:loghyp}
 \sup_j\int\log(1+F_{\epsilon_j})\,d\Pi_A\le C_0<\infty,
\end{equation}
then the entire selected joint law has a density \(f\) satisfying
\[
 \int f\log(1+f)\,d\nu\,dt\le C_0.
\]
In particular it gives positive-length distance sets for a positive
\(\nu\)-measure family of pins. No fixed exponent greater than one
is required by this hypothesis.
\end{proposition}

\begin{proof}
Markov's inequality bounds
\(\Pi_A\{F_{\epsilon_j}>T\}\) by \(C_0/\log(1+T)\).
Proposition~\ref{crowding:mass} first implies full absolute continuity.
To obtain the stronger entropy conclusion, let
\(f_j(y,t)=b_{\epsilon_j}*\sigma_y(t)\).
Whenever \(|t-|x-y||\le\epsilon_j/2\), the interval defining
\(f_j(y,t)\) is contained in
\([|x-y|-\epsilon_j,|x-y|+\epsilon_j]\), so
\(f_j(y,t)\le H_{\epsilon_j}(x,y)\).
Tonelli yields
\begin{align*}
 \int f_j\log(1+f_j)\,d\nu\,dt
 &\le\int\log(1+H_{\epsilon_j})\,d\Pi_A\\
 &\le\int\log(1+F_{\epsilon_j})\,d\Pi_A
       +\int N_{\epsilon_j}\,d\Pi_A\\
 &\le C_0+C\epsilon_j^\kappa.
\end{align*}
Here \(\log(1+F+N)\le\log(1+F)+N\).
Fibrewise differentiation of the now absolutely continuous law gives
\(f_j\to f\) for \(d\nu\,dt\)-almost every point.
Fatou proves the claimed entropy bound.
\end{proof}

The logarithmic test is weaker than a uniform positive fractional moment
of \(F_\epsilon\). The distinction is real even for ordinary scalar
densities: the probability density
\[
 f(t)=\frac{2\,\mathbf1_{(0,e^{-1})}(t)}
              {t(\log(1/t))^3}
\]
belongs to \(L\log L\) but to no \(L^{1+\beta}\), \(\beta>0\).
For completeness, its own crowding averages have uniformly bounded
logarithmic moments. Put
\(H_\epsilon(t)=\epsilon^{-1}\int_{t-\epsilon}^{t+\epsilon}f(u)du\),
and let \(I\) be a fixed finite interval containing the support.
Jensen applied with probability \(f(t)dt\) gives
\[
 \int f\log(1+H_\epsilon)
 \le\int f\log f+\log\!\int_I(1+H_\epsilon)
 \le\int f\log f+\log(|I|+2).
\]
The substitution \(v=\log(1/t)\) verifies both the entropy integrability
and the failure of every higher power. This scalar comparison asserts no
new planar Frostman construction or dimension theorem.

The substantive missing estimate is now explicit: beyond the curve
\eqref{hardgap:curve}, prove positive mass below some fixed far-crowding
threshold, or prove one of the stronger tail or moment bounds above.
The dimensions alone have not supplied any of these new bounds here.
